% Capítulo 4 - Desarrollo, pruebas y resultados.
% Es el capítulo más extenso. Fuentes: docs/requirements.md,
% docs/architecture.md, docs/diagrams.md, claude/decisiones.md
% (DEC-001 a DEC-036), código en src/, tests en tests/.
% Extensión objetivo: 18-22 páginas.

\chapter{Desarrollo, pruebas y resultados}
\label{ch:desarrollo}

% TODO: introducción del capítulo, ~media página.
% "Este capítulo describe el ciclo completo de desarrollo del PoC,
% organizado en seis bloques: análisis de requisitos, diseño
% arquitectónico, implementación, estrategia de pruebas, demostración
% funcional y comparativa cualitativa con los trabajos previos del
% tutor. El análisis detallado de requisitos se relega al apéndice
% \ref{ap:requirements}; aquí se presenta el resumen estructurado."

\section{Análisis de requisitos}
\label{sec:analisis}

% TODO: redactar. Resumen de docs/requirements.md (versión completa en
% apéndice B). Estructura:
% 4.1.1 Actores y casos de uso principales
%   - Tabla de actores con roles del orquestador y de la capa docente.
%   - Casos de uso del núcleo y de la capa docente, los más
%     representativos (no todos).
% 4.1.2 Requisitos funcionales (agrupados por área temática, no la
%   lista exhaustiva).
% 4.1.3 Requisitos no funcionales (tabla).
% 4.1.4 Restricciones y fuera de alcance.
%
% Texto típico: "El análisis completo se desarrolla en el apéndice
% \ref{ap:requirements}; este apartado presenta los elementos
% estructurantes."

\section{Diseño arquitectónico}
\label{sec:diseno}

% TODO: redactar. Insumo: docs/architecture.md + docs/diagrams.md.
% Importante: refrescar ambos documentos antes de escribir esta sección
% (les faltan Sprint 2 Pieza 3 y todo Sprint 3).
%
% Estructura:
% 4.2.1 Arquitectura en dos capas (DEC-004)
%   - Diagrama de contexto C4 (figura).
%   - Diagrama de contenedores C4 (figura).
% 4.2.2 Modelo del dominio
%   - Árbol de descriptores con materialized path (DEC-008).
%   - Herencia en cascada por tipo de dato (DEC-026).
%   - Cierre léxico de plantillas (refinamiento DEC-012, 2025-06-06).
% 4.2.3 Patrones aplicados
%   - Facade (OrchestratorService, DEC-031).
%   - Command pattern (DEC-028).
%   - Repository (DEC-007/030).
%   - Strategy / ABC para conectores (DEC-029).
%   - Protocol para la frontera CLI/cliente HTTP (DEC ...).
% 4.2.4 Frontera REST como API pública
%   - Argumento defendible: la API REST no se monta para que la CLI
%     remota funcione (el admin tiene SSH), sino para que existan
%     integraciones externas.
%   - Cliente CLI dual mode in-process / HTTP descartado a favor del
%     modelo cliente puro (Sprint 3.2 B-2), por coherencia con docker /
%     kubectl.
% 4.2.5 Diagrama de componentes (figura) y de despliegue (figura).

\section{Implementación}
\label{sec:implementacion}

% TODO: redactar. Estructura por capas, no por sprints (los sprints
% importan para la metodología, no para la presentación del producto).
% 4.3.1 Núcleo del orquestador
%   - Modelos Pydantic.
%   - Resolver puro (función + tests con Hypothesis).
%   - Motor de Jobs y Commands.
%   - Repositorios sobre TinyDB.
%   - Servicio facade.
% 4.3.2 Conector mock
%   - Estado en memoria, plantillas demo precargadas.
%   - Sin dependencias externas: arranca el sistema entero sin
%     hipervisor real.
% 4.3.3 Interfaz CLI
%   - cmd2 con argparse, modo doble REPL/bash.
%   - Cliente HTTP por inyección, traduce errores de transporte a
%     OperationError con mensaje accionable.
% 4.3.4 Daemon REST
%   - FastAPI + uvicorn, paths como query params.
%   - Healthcheck, codificación de errores 200/400/422.
% 4.3.5 Empaquetado y distribución
%   - Dockerfile multi-stage, docker-compose con volumen named.
%   - Instalador install.sh con detección de Docker, puerto, Compose.
%   - Imagen publicada en ghcr.io con tags semver.

\section{Estrategia de pruebas}
\label{sec:pruebas}

% TODO: redactar.
% 4.4.1 Pirámide de pruebas aplicada.
% 4.4.2 Unidad: por módulo, con cobertura alta en lógica de dominio.
% 4.4.3 Integración: con TinyDB en memoria, mock connector real.
% 4.4.4 Property-based con Hypothesis: el resolver, función pura, es
%   el caso ideal. Tres propiedades comprobadas.
% 4.4.5 REST: FastAPI TestClient como cliente httpx que enruta al ASGI
%   in-process. Caza divergencias cliente-servidor sin red.
% 4.4.6 Métricas finales: 244 tests verdes al cierre de Sprint 3,
%   95-100% en módulos centrales.

\section{Demostración funcional}
\label{sec:demo}

% TODO: redactar. Capturas a incluir (carpeta figures/).
% - Estructura del árbol cargada desde el demo.
% - inspect mostrando herencia y closure.
% - Ciclo completo deploy/start/status/stop/undeploy.
% - Comprobación HTTP desde curl y desde el navegador (/docs).
% - Recorrido por la API en Swagger UI.

\section{Comparativa con esxobjects y yamlinfr}
\label{sec:comparativa}

% TODO: redactar.
% - Tabla comparativa: alcance, dependencias, modelo de descriptores,
%   capacidad de extensión, cobertura de pruebas, distribución.
% - Lectura cualitativa: yamlinfr/esxobjects resolvían su caso, atado a
%   VMware. UnivOrch generaliza y abre la puerta a otros hipervisores
%   sin tocar el núcleo.
% - Discusión honesta: UnivOrch no es un sustituto operativo todavía
%   (le falta el conector real); es un PoC del modelo.
